% =============================================================================
% CompetitiveExam/CS301/07-linked-lists-doubly-applications-questions.tex
% Competitive Exam Practice 7 — GATE (Computer Science) Pattern
% Doubly linked lists: insert/delete pointer surgery and the two-link invariant,
% O(1) delete-by-pointer vs singly O(n), circular singly/doubly lists and
% NULL-free traversal termination, list-based stack/queue (linked circular
% queue), polynomial representation and merge-by-descending-exponent addition.
%
% Student-facing — NO answers. Answer key: 07-linked-lists-doubly-applications-answers.tex.
%
% Provenance (availability-first, Tier 1 = local PYQ book
% CompetitiveExam/Books/GATE-PYQs/filter1_volume2.pdf):
%   - Q1  [GATE CSE 1997 Q1.4, verified — GATEOverflow V2 p.233] —
%     O(1) concatenation → circular doubly linked list. Stem + all four
%     options confirmed in the local text layer (volume2, Linked List cluster,
%     title 3.10.6) and independently on PracticePaper's GATE CSE 1997 paper
%     page (same stem + options A-D). Local book answer key: 3.10.6 = C.
%     Answer (C).
%   - Q2  [GATE CSE 2004 Q36, verified — GATEOverflow V2 p.235] —
%     circular list as a queue, one variable to access it → point at the rear
%     node. Stem + all four options in the local text layer (title 3.10.12);
%     local book answer key: 3.10.12 = A. Answer (A).
%   - Q3  [GATE CSE 2017 Set 2 Q13, verified — GATEOverflow V2 p.242] —
%     circular queue on a singly linked list with FRONT/REAR: rear->next must
%     point to front. Stem + options (I)/(II) + A-D in the local text layer
%     (title 3.12.10); local book answer key: 3.12.10 = B. Answer (B).
%   - Q4  [GATE CSE 2018 Q3, verified — GATEOverflow V2 p.242] —
%     non-circular singly linked list queue, enqueue at head / dequeue at tail
%     → Θ(1), Θ(n). Stem in the local text layer (title 3.12.11); all four
%     options recovered from GeeksforGeeks' GATE-CS-2018 linked-list question
%     page (Θ(1),Θ(1) / Θ(1),Θ(n) / Θ(n),Θ(1) / Θ(n),Θ(n)); local book answer
%     key: 3.12.11 = B. Answer (B).
%   - Q5  [GATE CSE 2023 Q3, verified — GATEOverflow V2 p.237] —
%     SLLdel (singly, given node pointer) is O(n), DLLdel (doubly) is O(1).
%     Stem in the local text layer (title 3.10.20); all four options recovered
%     from GeeksforGeeks' GATE-CS-2023 question page; answer D corroborated by
%     the GATEQA listing ("Answer: D") and the local book answer key
%     (3.10.20 = D). Answer (D).
%   - Q6–Q10  [Original, GATE-pattern] — fresh instances (the deck, Notes, and
%     the Week 7 exercises use 10-20-30, 5-15-25, a-b-c-d, 4x^4+2x^2+7, and
%     30-20-10/40; none of those concrete values is reused here). Every trace
%     was checked by hand this session:
%       Q6  DLL delete-by-pointer on head <-> 7 <-> 19 <-> 42: removes 19,
%           rewrites 7->next and 42->prev in O(1) once p is known. Verified.
%       Q7  DLL insert-after on head <-> 3 <-> 8 <-> 14, insert 5 after 8:
%           the four-write order that restores the invariant. Verified.
%       Q8  singly circular head -> p -> q -> r -> p under while (ptr != NULL):
%           never meets NULL, cycles forever. Verified.
%       Q9  merge-add A = 6x^3 + 2x + 9 and B = 4x^3 + 5x^2 - 2x:
%           10x^3 + 5x^2 + 9 (the x term cancels), one pass O(n+m). Verified.
%       Q10 list stack head -> 7 -> 14 -> 21: push(28) then pop() -> 28, top 7,
%           both O(1) worst case. Verified.
%
% Method note: no Ghostscript/Tesseract OCR in this environment, so
% image-embedded math/options could not be OCR'd; they were recovered instead
% by (a) the local pdftotext text layer where present and (b) direct web
% fetches of GeeksforGeeks/PracticePaper question pages (search engines were
% partly bot-blocked; no forked claim-verifier Task tool exists here, so CoVe
% was executed inline against the local answer-key section + the web pages).
% The local book's answer-key section — which the registry assumed was
% collapsed — is in fact recoverable via pdftotext and was used for every
% real-PYQ answer.
%
% Compile with: cd CompetitiveExam/CS301 &&
%   TEXINPUTS=../../Preambles;$TEXINPUTS xelatex -interaction=nonstopmode 07-linked-lists-doubly-applications-questions.tex
%   (Windows/MiKTeX uses ; not : as the separator)
% =============================================================================

\documentclass[11pt]{article}
\usepackage[margin=1in]{geometry}
% =============================================================================
% Preambles/header.tex — shared LaTeX/Beamer preamble for this project
%
% Use from a Beamer lecture by prepending:
%     \input{header}
% and compiling with TEXINPUTS=../Preambles:$TEXINPUTS (the /compile-latex
% skill does this automatically).
%
% PALETTE CONTRACT. Color names defined here MUST match the names in
% Quarto/theme-template.scss so Beamer and Quarto renderings use the same
% palette. The scripts/check-palette-sync.sh script enforces this.
%
% To customize: edit the HEX values below AND the matching $variable in
% Quarto/theme-template.scss, then run scripts/check-palette-sync.sh.
% =============================================================================

% -----------------------------------------------------------------------------
% Engine detection + encoding
%
% Our /compile-latex skill uses XeLaTeX, where inputenc is unnecessary and
% can produce encoding warnings. Only load inputenc under pdfTeX.
% -----------------------------------------------------------------------------
\usepackage{iftex}
\ifPDFTeX
  \usepackage[utf8]{inputenc}
\fi

\usepackage{amsmath, amssymb, amsthm}
\usepackage{booktabs}
\usepackage{graphicx}

% -----------------------------------------------------------------------------
% PALETTE — mirrors Quarto/theme-template.scss variable names.
% Load xcolor and define colors BEFORE anything that references them
% (notably hyperref, hypersetup, and the Beamer color assignments).
% -----------------------------------------------------------------------------
\usepackage{xcolor}

\definecolor{primary-blue}{HTML}{012169}      % headings, accents
\definecolor{primary-gold}{HTML}{B9975B}      % emphasis, borders
\definecolor{highlight-yellow}{HTML}{F2A900}  % markers, alerts
\definecolor{light-bg}{HTML}{E8EDF5}          % title / section backgrounds
\definecolor{jet}{HTML}{1A1A1A}               % body text

% Semantic colors (used in diagrams and annotations)
\definecolor{positive}{HTML}{15803D}          % good / observed / identified
\definecolor{negative}{HTML}{B91C1C}          % bad / problematic / bias
\definecolor{neutral}{HTML}{525252}           % reference / context

% Highlight accents (match .hi-* classes in the SCSS)
\definecolor{hi-slate}{HTML}{314F4F}
\definecolor{hi-green}{HTML}{2E8B57}
\definecolor{hi-red}{HTML}{C41E3A}

% Semantic aliases so skill templates and snippets can write
% \textcolor{accent}{...} without hard-coding "primary-blue".
\colorlet{accent}{primary-blue}
\colorlet{accent2}{primary-gold}

% -----------------------------------------------------------------------------
% Hyperref — loaded AFTER xcolor + palette so link colors resolve.
% -----------------------------------------------------------------------------
\usepackage{hyperref}
\hypersetup{colorlinks=true, linkcolor=primary-blue, urlcolor=primary-blue,
            citecolor=primary-blue, pdfborder={0 0 0}}

% -----------------------------------------------------------------------------
% Course code — set once per deck (after \input{header}, before \begin{document})
% via \coursecode{CS401}. Shown in the footer on every slide so a deck's course
% is identifiable without opening the title page. Empty by default.
% -----------------------------------------------------------------------------
\makeatletter
\newcommand{\coursecode}[1]{\gdef\@coursecodetext{#1}}
\gdef\@coursecodetext{}
\makeatother

% -----------------------------------------------------------------------------
% Beamer theme assignments (only apply under Beamer).
%
% We detect Beamer via \@ifclassloaded{beamer}, which is the documented,
% non-internal way. This must be wrapped in \makeatletter/\makeatother
% because \@ifclassloaded contains an @-catcode character.
% -----------------------------------------------------------------------------
\makeatletter
\@ifclassloaded{beamer}{%
  \usefonttheme{professionalfonts}
  \setbeamercolor{structure}{fg=primary-blue}
  \setbeamercolor{frametitle}{fg=primary-blue}
  \setbeamercolor{title}{fg=primary-blue}
  \setbeamercolor{subtitle}{fg=primary-gold}
  \setbeamercolor{author}{fg=primary-blue}
  \setbeamercolor{institute}{fg=neutral}
  \setbeamercolor{date}{fg=primary-gold}
  \setbeamercolor{itemize item}{fg=primary-blue}
  \setbeamercolor{itemize subitem}{fg=primary-gold}
  \setbeamercolor{alerted text}{fg=primary-gold}
  \setbeamercolor{block title}{fg=white, bg=primary-blue}
  \setbeamercolor{block body}{bg=light-bg}
  % Semantic block variants: block=definition/concept (blue), exampleblock=
  % worked example (green/positive), alertblock=key takeaway or caution (gold).
  % Keep to <=2 colored boxes per slide across ALL three variants (INV-7).
  \setbeamercolor{block title example}{fg=white, bg=positive}
  \setbeamercolor{block body example}{bg=light-bg}
  \setbeamercolor{block title alerted}{fg=white, bg=primary-gold}
  \setbeamercolor{block body alerted}{bg=light-bg}

  % Minimal footer: course code (left) + page X / N (right), no navigation clutter
  \setbeamertemplate{navigation symbols}{}
  \setbeamertemplate{footline}{%
    \color{neutral}\footnotesize\hspace{0.7em}\@coursecodetext\hfill
    \insertframenumber/\inserttotalframenumber\hspace{0.7em}\vspace{0.3em}}
}{}
\makeatother

% -----------------------------------------------------------------------------
% TikZ setup — libraries the snippet gallery uses
% -----------------------------------------------------------------------------
\usepackage{tikz}
\usetikzlibrary{arrows.meta, positioning, calc, decorations.pathreplacing,
                fit, shapes.geometric, backgrounds}

% Shared TikZ styles — snippets and hand-written diagrams can pick these up
% via `\begin{tikzpicture}[every node/.style={...}]` or by naming specific
% styles (e.g., `\node[dag-node] (X) at (0,0) {X};`).
\tikzset{
  % Node styles for causal DAGs and similar
  dag-node/.style = {
    draw=primary-blue, fill=light-bg, rounded corners=2pt,
    minimum width=1.2cm, minimum height=0.9cm, align=center, font=\small
  },
  decision-node/.style = {
    draw=primary-gold, fill=white, diamond, aspect=1.8,
    minimum width=1.8cm, minimum height=1.2cm, align=center, font=\small,
    inner sep=1pt
  },
  % Edge styles — observed vs counterfactual
  observed-edge/.style = {-{Latex[length=2.5mm]}, thick, draw=jet},
  counterfactual-edge/.style = {-{Latex[length=2.5mm]}, thick, draw=neutral, dashed},
  confound-edge/.style = {-{Latex[length=2.5mm]}, thick, draw=negative, dashed},
  % Data-point styles for plots
  observed-dot/.style = {fill=primary-blue, draw=primary-blue, circle, inner sep=1.2pt},
  counterfactual-dot/.style = {fill=white, draw=neutral, circle, inner sep=1.2pt}
}

% -----------------------------------------------------------------------------
% Convenience macros
% -----------------------------------------------------------------------------
\newcommand{\muted}[1]{\textcolor{neutral}{#1}}
\newcommand{\key}[1]{\textcolor{primary-gold}{\textbf{#1}}}
\newcommand{\good}[1]{\textcolor{positive}{#1}}
\newcommand{\bad}[1]{\textcolor{negative}{#1}}

% Standout frame macro: full-bleed dark background for section transitions.
% Beamer-only (uses \begin{frame}/\setbeamercolor/\usebeamerfont) -- do not
% call from an article-class file (e.g. Notes/<CODE>/*-notes.tex). \muted,
% \key, \good, \bad, and \coursecode above are plain xcolor/gdef and are
% safe to use from either class; verified when Notes/article-class support
% was added.
% Expand: \transitionslide{Your transition title}
\newcommand{\transitionslide}[1]{%
  \begingroup
  \setbeamercolor{background canvas}{bg=primary-blue}
  \begin{frame}[plain]
    \centering
    \vfill
    {\usebeamerfont{title}\color{white}\Large #1\par}
    \vfill
  \end{frame}
  \endgroup
}

% Section divider frame (Beamer-only), an alternative to \transitionslide with a
% darker two-line style: near-black (jet) background, a small white label line
% above a larger gold title, and a thin gold rule along the bottom edge.
% Expand: \sectiondivider{Section 2}{Pointers}
\newcommand{\sectiondivider}[2]{%
  \begingroup
  \setbeamercolor{background canvas}{bg=jet}
  \begin{frame}[plain]
    \vspace*{3.0cm}
    {\color{white}\ttfamily\large #1\par}
    \vspace{0.5cm}
    {\color{primary-gold}\LARGE #2\par}
    \begin{tikzpicture}[remember picture, overlay]
      \draw[primary-gold, line width=2.5pt]
        ($(current page.south west)+(0,0.1)$) -- ($(current page.south east)+(0,0.1)$);
    \end{tikzpicture}
  \end{frame}
  \endgroup
}

% -----------------------------------------------------------------------------
% Channel watermark — "Concepts That Click" (YouTube).
%
% Article-class files (CompetitiveExam Q/A sets, Notes, Assignments) get a
% light diagonal draft watermark on every page plus a centered footer naming
% the channel. Beamer decks get a subtle TikZ background watermark instead
% (the footline already carries the course code). Centralized here so every
% MCQ PDF is branded without per-file edits.
% -----------------------------------------------------------------------------
\makeatletter
\@ifclassloaded{article}{%
  \usepackage{draftwatermark}
  \SetWatermarkText{Concepts That Click}
  \SetWatermarkScale{1}
  \SetWatermarkFontSize{2cm}
  \SetWatermarkLightness{0.86}
  \SetWatermarkAngle{45}
  \usepackage{fancyhdr}
  \pagestyle{fancy}
  \fancyhf{}
  \fancyfoot[C]{\footnotesize\textcolor{neutral}{Concepts That Click~|~YouTube: \href{https://www.youtube.com/@concepts.thatclick}{@concepts.thatclick}}}
  \renewcommand{\headrulewidth}{0pt}
  \renewcommand{\footrulewidth}{0pt}
  \fancypagestyle{plain}{%
    \fancyhf{}%
    \fancyfoot[C]{\footnotesize\textcolor{neutral}{Concepts That Click~|~YouTube: \href{https://www.youtube.com/@concepts.thatclick}{@concepts.thatclick}}}%
    \renewcommand{\headrulewidth}{0pt}%
    \renewcommand{\footrulewidth}{0pt}%
  }
}{}
\@ifclassloaded{beamer}{%
  \setbeamertemplate{background}{%
    \begin{tikzpicture}[remember picture, overlay]
      \node[rotate=45, opacity=0.055, align=center]
        at (current page.center)
        {\Huge\textcolor{neutral}{Concepts That Click}};
    \end{tikzpicture}%
  }
}{}
\makeatother 
\coursecode{CS301}
\renewcommand{\thesection}{7.\arabic{section}}

\title{Competitive Exam Practice 7: Doubly \& Circular Lists, List-Backed Stack and Queue, and Polynomials\\\large GATE (Computer Science) Pattern}
\author{Auqib Hamid Lone\\[0.3em]
  \emph{Department of Computer Science \& Engineering}, \emph{GCET Ganderbal}}
\date{\today}

\begin{document}
\maketitle

\noindent Each question carries 1--2 marks. MCQ = single correct option; NAT = Numerical Answer Type.

\begin{enumerate}
\item[\textbf{Q1}] \textit{[GATE CSE 1997 Q1.4, verified --- GATEOverflow V2 p.233]} \textbf{[MCQ, 1 mark]} The concatenation of two lists is to be performed in $O(1)$ time. Which of the following implementations of a list should be used?
\begin{enumerate}
  \item[(A)] Singly linked list
  \item[(B)] Doubly linked list
  \item[(C)] Circular doubly linked list
  \item[(D)] Array implementation of list
\end{enumerate}

\item[\textbf{Q2}] \textit{[GATE CSE 2004 Q36, verified --- GATEOverflow V2 p.235]} \textbf{[MCQ, 1 mark]} A circularly linked list is used to represent a queue. A single variable \texttt{p} is used to access the queue. To which node should \texttt{p} point such that both the operations \texttt{enqueue} and \texttt{dequeue} can be performed in constant time?
\begin{enumerate}
  \item[(A)] rear node
  \item[(B)] front node
  \item[(C)] not possible with a single pointer
  \item[(D)] node next to front
\end{enumerate}

\item[\textbf{Q3}] \textit{[GATE CSE 2017 Set 2 Q13, verified --- GATEOverflow V2 p.242]} \textbf{[MCQ, 2 marks]} A circular queue has been implemented using a singly linked list where each node consists of a value and a single pointer pointing to the next node. We maintain exactly two external pointers \texttt{FRONT} and \texttt{REAR} pointing to the front node and the rear node of the queue, respectively. Which of the following statements is/are CORRECT for such a circular queue, so that insertion and deletion operations can be performed in $O(1)$ time?
\begin{enumerate}
  \item[I.] Next pointer of front node points to the rear node.
  \item[II.] Next pointer of rear node points to the front node.
\end{enumerate}
\begin{enumerate}
  \item[(A)] (I) only
  \item[(B)] (II) only
  \item[(C)] Both (I) and (II)
  \item[(D)] Neither (I) nor (II)
\end{enumerate}

\item[\textbf{Q4}] \textit{[GATE CSE 2018 Q3, verified --- GATEOverflow V2 p.242]} \textbf{[MCQ, 1 mark]} A queue is implemented using a non-circular singly linked list. The queue has a head pointer and a tail pointer. Let $n$ denote the number of nodes in the queue. Let \texttt{enqueue} be implemented by inserting a new node at the head, and \texttt{dequeue} be implemented by deletion of a node from the tail. Which one of the following is the time complexity of the most time-efficient implementation of \texttt{enqueue} and \texttt{dequeue}, respectively, for this data structure?
\begin{enumerate}
  \item[(A)] $\Theta(1)$, $\Theta(1)$
  \item[(B)] $\Theta(1)$, $\Theta(n)$
  \item[(C)] $\Theta(n)$, $\Theta(1)$
  \item[(D)] $\Theta(n)$, $\Theta(n)$
\end{enumerate}

\item[\textbf{Q5}] \textit{[GATE CSE 2023 Q3, verified --- GATEOverflow V2 p.237]} \textbf{[MCQ, 1 mark]} Let \texttt{SLLdel} be a function that deletes a node in a singly linked list given a pointer to the node and a pointer to the head of the list. Similarly, let \texttt{DLLdel} be another function that deletes a node in a doubly linked list given a pointer to the node and a pointer to the head of the list. Let $n$ denote the number of nodes in each of the linked lists. Which one of the following choices is TRUE about the worst-case time complexity of \texttt{SLLdel} and \texttt{DLLdel}?
\begin{enumerate}
  \item[(A)] \texttt{SLLdel} is $O(1)$ and \texttt{DLLdel} is $O(n)$
  \item[(B)] Both \texttt{SLLdel} and \texttt{DLLdel} are $O(\log n)$
  \item[(C)] Both \texttt{SLLdel} and \texttt{DLLdel} are $O(n)$
  \item[(D)] \texttt{SLLdel} is $O(n)$ and \texttt{DLLdel} is $O(1)$
\end{enumerate}

\item[\textbf{Q6}] \textit{[Original, GATE-pattern]} \textbf{[MCQ, 2 marks]} A doubly linked list holds \texttt{head} $\leftrightarrow$ 7 $\leftrightarrow$ 19 $\leftrightarrow$ 42 $\leftrightarrow$ \texttt{NULL}. A pointer \texttt{p} points at the node holding~19. That node is deleted by the deck's delete-by-pointer and then \texttt{free(p)} is called. Which statement is TRUE about the resulting list and the cost?
\begin{enumerate}
  \item[(A)] The list becomes \texttt{head} $\leftrightarrow$ 7 $\leftrightarrow$ 42; 7's \texttt{next} is set to 42 and 42's \texttt{prev} is set to 7, both reached in $O(1)$ from \texttt{p}; the delete is $O(1)$ worst case once \texttt{p} is known.
  \item[(B)] The list becomes \texttt{head} $\leftrightarrow$ 7 $\leftrightarrow$ 42, but the predecessor of the deleted node must first be found by a walk, so the delete is $O(n)$ worst case.
  \item[(C)] The list is unchanged --- \texttt{head} $\leftrightarrow$ 7 $\leftrightarrow$ 19 $\leftrightarrow$ 42 --- because delete-by-pointer only moves \texttt{head}.
  \item[(D)] The list becomes \texttt{head} $\leftrightarrow$ 7 $\leftrightarrow$ 42, but only 7's \texttt{next} is updated; 42's \texttt{prev} still points at the freed node, breaking the two-link invariant.
\end{enumerate}

\item[\textbf{Q7}] \textit{[Original, GATE-pattern]} \textbf{[MCQ, 2 marks]} A doubly linked list holds \texttt{head} $\leftrightarrow$ 3 $\leftrightarrow$ 8 $\leftrightarrow$ 14 $\leftrightarrow$ \texttt{NULL}, and \texttt{p} points at the node holding~8. A new node holding~5 is inserted immediately after \texttt{p}. Which sequence of pointer writes restores the two-link invariant and yields \texttt{head} $\leftrightarrow$ 3 $\leftrightarrow$ 8 $\leftrightarrow$ 5 $\leftrightarrow$ 14?
\begin{enumerate}
  \item[(A)] \texttt{newNode->next = p->next;}\\
  \texttt{newNode->prev = p;}\\
  \texttt{if (p->next != NULL) p->next->prev = newNode;}\\
  \texttt{p->next = newNode;}
  \item[(B)] \texttt{p->next = newNode;}\\
  \texttt{newNode->prev = p;}\\
  \texttt{newNode->next = p->next;}
  \item[(C)] \texttt{newNode->next = p->next;}\\
  \texttt{p->next = newNode;}\\
  \texttt{newNode->prev = p;}\\
  \texttt{if (p->next != NULL) p->next->prev = newNode;}
  \item[(D)] \texttt{newNode->prev = p;}\\
  \texttt{newNode->next = p->next;}\\
  \texttt{p->next = newNode;}\\
  \texttt{p->next->prev = newNode;}
\end{enumerate}

\item[\textbf{Q8}] \textit{[Original, GATE-pattern]} \textbf{[MCQ, 1 mark]} A singly circular list reads \texttt{head} $\to$ p $\to$ q $\to$ r $\to$ p. The following loop is run:
\begin{verbatim}
ptr = head;
while (ptr != NULL) {
    visit(ptr);
    ptr = ptr->next;
}
\end{verbatim}
What happens?
\begin{enumerate}
  \item[(A)] It visits p, q, r and then stops when \texttt{ptr} becomes \texttt{NULL}.
  \item[(B)] It loops forever, cycling p, q, r, p, q, r, $\dots$ and never meets \texttt{NULL}.
  \item[(C)] It visits only the node p, then stops at the next \texttt{head}.
  \item[(D)] It aborts after three visits because the list has no end.
\end{enumerate}

\item[\textbf{Q9}] \textit{[Original, GATE-pattern]} \textbf{[MCQ, 2 marks]} Two sparse polynomials are stored as ordered $(\text{coeff},\text{exp})$ chains in strictly descending exponent order: $A = 6x^{3} + 2x + 9$ and $B = 4x^{3} + 5x^{2} - 2x$. Using the deck's merge-by-descending-exponent addition, which result and cost are correct?
\begin{enumerate}
  \item[(A)] $A + B = 10x^{3} + 5x^{2} + 9$, produced by one merge pass costing $O(n+m)$ worst case; the cancelled $x$ term is never stored.
  \item[(B)] $A + B = 10x^{3} + 5x^{2} + 0x + 9$; the merge stores the zero term, still costing $O(n+m)$.
  \item[(C)] $A + B = 10x^{3} + 5x^{2} + 9$, but the merge costs $O(n \cdot m)$ because every pair of terms is multiplied.
  \item[(D)] $A + B = 10x^{3} + 5x^{2} + 4x + 9$, produced by one merge pass costing $O(n+m)$.
\end{enumerate}

\item[\textbf{Q10}] \textit{[Original, GATE-pattern]} \textbf{[MCQ, 1 mark]} A linked-list stack has \texttt{head} $\to$ 7 $\to$ 14 $\to$ 21 $\to$ \texttt{NULL}, with the top of the stack at \texttt{head}. \texttt{push(28)} is performed, then one \texttt{pop()}. Which statement is TRUE?
\begin{enumerate}
  \item[(A)] \texttt{push} links 28 in front, giving \texttt{head} $\to$ 28 $\to$ 7 $\to$ 14 $\to$ 21; \texttt{pop()} returns 28 and the new top is 7; both operations are $O(1)$ worst case.
  \item[(B)] \texttt{push} appends 28 at the tail, giving \texttt{head} $\to$ 7 $\to$ 14 $\to$ 21 $\to$ 28; \texttt{pop()} removes the front, so it returns 7.
  \item[(C)] \texttt{push} links 28 in front, giving \texttt{head} $\to$ 28 $\to$ 7 $\to$ 14 $\to$ 21; \texttt{pop()} returns 7 because a stack removes the oldest item.
  \item[(D)] Both \texttt{push} and \texttt{pop} take $O(n)$ worst case, because a linked list has no indices.
\end{enumerate}

\end{enumerate}

\end{document}
